\section{Preliminaries and Overview} \label{sec:overview}

\subsection{Shape Representation}
We represent the target shape $\target$ and each object $\object_i$ as unions of smooth convex bodies modeled by differentiable support functions (DSFs)~\cite{lee2023uncertain}.
For a compact convex body $\convex\subset\real^3$, its support function is $\dsf_\convex(x)=\max_{p\in\convex}x^\top p$ for $x\ne0$.
At pose $q=(R,t)\in\SE(3)$, we write
\begin{equation*}
    \convex(q)=\{Rp+t\mid p\in\convex\},
    \qquad \object_i(q_i)=\bigcup_j\convex_{ij}(q_i),
\end{equation*}
where $\convex_{ij}$ are the convex components in the local frame of object $i$.
Similarly, $\target=\bigcup_j\target_j$, where $\target_j$ are fixed convex components in the world frame.
DSFs support differentiable signed-distance and contact-geometry computations between posed convex bodies~\cite{an2024riemannianmtd}, enabling gradient-based pose optimization.

% Queue the pose-optimization figure on page 3 after the overview introduction.

\subsection{Markov Decision Process Formulation}

Given an inventory of $n$ objects, we model stacking as a finite-horizon sequential decision process in which one object is placed at each step.
Let $\placed_k\subseteq\{1,\ldots,n\}$ be the indices of the objects already placed, and let $i_0\notin\{1,\ldots,n\}$ index the fixed ground.
The current scene is $\scene_k=\{(i,q_i)\mid i\in\placed_k\}$, and the remaining inventory is $\inventory_k=\{1,\ldots,n\}\setminus\placed_k$.
The state, action, and deterministic transition function are
\begin{equation*}
    \state_k=(\scene_k,\inventory_k,\target),\quad
    \action_k=(i_{k+1},q_{k+1}),\quad
    \state_{k+1}=\transition(\state_k,\action_k),
\end{equation*}
where $i_{k+1}\in\inventory_k$. The transition inserts the proposed pose into the scene, simulates the resulting structure, and returns the settled scene with $\inventory_{k+1}=\inventory_k\setminus\{i_{k+1}\}$. 
The reward $\reward(\state_k,\action_k)$ scores volume filling, equilibrium residuals, and disturbance response; feasibility checks separately reject invalid transitions (Sec.~\ref{sec:sequence_planning}).
The MDP is $\mdp=(\mathcal{S},\mathcal{A},\transition,\reward,\gamma)$, with discount factor $\gamma$.

\subsection{MCTS-Based Stacking Framework}

Our framework uses MCTS to search placement sequences, integrating pose sampling, bilevel refinement, and simulation within each expansion (Fig.~\ref{fig:overview}).
The search tree $\tree=(\nodes,\edges)$ is defined over the MDP.
Each node $\node\in\nodes$ represents a state $\state(\node)$, and each edge $(\node,a)\in\edges$ represents an action and its transition to a child node.
Let $\visit(\node)$ denote the number of completed evaluations and $Q(\node)$ the empirical mean of discounted returns backpropagated to that node.

Each iteration balances value and exploration to select a path, expands and evaluates an action, then backpropagates its discounted return to update $Q$ and $\visit$.

Expansion uses three modules (Fig.~\ref{fig:overview}): $\initpose$ proposes an object and initial pose $\hat{\action}_k=(i_{k+1},\hat q_{k+1})$, $\poseopt$ refines the pose, and $\simul$ determines feasibility, the settled successor state, and reward.
The resulting flow is
\begin{equation*}
\state_k
\xrightarrow{\ \initpose\ }
\hat{\action}_k
\xrightarrow{\ \poseopt\ }
\action_k
\xrightarrow{\ \simul\ }
\state_{k+1},\ \reward(\state_k,\action_k)
\end{equation*}
Independent candidates follow this flow in parallel on the GPU, with bodywise ADMM also parallelizing contact force optimization within each structure (Sec.~\ref{sec:pose_optimization}).
After search, we execute the selected placement and replan from the observed scene (Sec.~\ref{sec:sequence_planning}).
